\subsection{\problemblock{*Survey} Data Block}\label{sec:block-survey}

The purpose of a survey is to change one or more input parameters and determine the
effect on the trajectories. This is a tradeoff or sensitivity study. Trajectories are
computed for every value of each survey variable.

As described in \cref{sec:problem-file-format}, any numerical value in a problem
file can be varied by replacing it with a survey placeholder. For example,

\begin{lstlisting}[style=taosmanual]
*initial alt=surv-1 lat=0 long=0 wt=1200
\end{lstlisting}

sets initial altitude equal to survey 1. The keyword \texttt{surv-1} is a placeholder;
the actual value comes from the corresponding \texttt{*survey 1} block. Survey
numbers relate placeholders and blocks: \texttt{surv-1} corresponds to
\texttt{*survey 1}, \texttt{surv-2} to \texttt{*survey 2}, and so on.

The placeholders determine which numerical locations are modified. The survey
blocks determine the values assigned to them. In the preceding example there must
be a \texttt{*survey 1} block that supplies initial-altitude values. For example,

\begin{lstlisting}[style=taosmanual]
*survey 1 alt0 lo=40000 hi=80000 inc=20000
\end{lstlisting}

runs three cases with initial altitude equal to 40,000, 60,000, and 80,000~ft. TAOS
substitutes the survey value everywhere the matching \texttt{surv-n} keyword appears,
so several input values can be changed together by one survey.

\taossubhead{Survey Numbers and Names}

A block begins with the \problemblock{*survey} keyword and the survey number.
Numbers are normally sequential beginning with one. Any number of survey loops may
be defined. They form nested loops: the lowest-numbered survey is innermost and the
highest-numbered survey is outermost.

A survey-variable name follows the number. It must not duplicate an output-variable
name. The name is used in EGS database files and in summary output to identify the
quantity being surveyed. In the example above, the database contains trajectories
associated with \texttt{alt0=40000}, another set associated with
\texttt{alt0=60000}, and so forth.

\taossubhead{Low, High, and Increment Survey Values}

Survey values can be assigned in two ways. In the incremental form, \statevar{lo}
provides the starting value. A set of trajectories is always run first at that value. After
the case is complete, \statevar{inc} is added. If the increment is zero, only the low
case is run.

After incrementing, TAOS compares the value with \statevar{hi}. If it exceeds the
high value, the survey value is set exactly equal to the high value and that case is
marked as the last. This process continues until the high case has been calculated.

\taossubhead{Specific Survey Values}

The second method enters the desired values directly:

\begin{lstlisting}[style=taosmanual]
*survey 1 alt0 vals=30000,60000,75000
\end{lstlisting}

This runs three cases with initial altitude set to 30,000, 60,000, and 75,000~ft and is
useful when the values do not follow a regular increment.

Both methods may be combined:

\begin{lstlisting}[style=taosmanual]
*survey 1 alt0 lo=40000 hi=80000 inc=20000
  vals=35000,45000
\end{lstlisting}

The survey then contains 35,000, 40,000, 45,000, 60,000, and 80,000. Values from
\statevar{lo}, \statevar{hi}, and \statevar{inc} are combined with the explicit
\statevar{vals} entries. The extra values are inserted in the survey in order from lowest to highest
value, so the first trajectory set uses \texttt{alt0=35000}, the second
\texttt{alt0=40000}, and so on.
